\chapter{Functional Inferences}

So far we have talked about the basics of Probability, and introduced key concepts like probability mass functions. While these abstractions are extremely useful and will follow us through the course of our lives as physicists, realistically we will rarely measure an observable extracted from a completely known pdf. Actually, if the probability density function is completely known, measuring the observable brings us no new information other than the measure itself - we learn nothing new about the world. 

Things start to get interesting when the distribution itself depends on some unknown parameter\sidenote{This parameter can be anything, it does not have to be observable, follow a pdf, be a random variable, it can be anything that the density function of the observable depends on.}, and measuring the observable brings us new information on what that parameter may be. In this section we will write
	\[
	p_X(x;\theta)
	\]
to denote the pdf of observable $x\in X$ dependent on some parameter $\theta\in M$, using the semicolon to distinguish observables from parameters. Our goal is that of the \emph{functional inference}, finding a \emph{function} that depends on the observables (statistic) that assigns a sort of "score" to every value of the parameters, indicating their plausibility. In short:
	\[
	s(x): x\rightarrow f_x(\theta); \quad f:A\rightarrow \R
	\]

\section{The Likelihood function}

We begin with one of the most important definitions of the whole course:

\begin{definition}[Likelihood function]
	Let $x\sim p(x;\theta)\, , \theta\in M$. Given an observation $x_0$ of $x$, we define \textbf{Likelihood function} the value of the probability (or density) of $x_0$ as a function of $\theta$
	\[
	L: M \rightarrow \R_0^+:L_{x_0}(\theta)=p_(x_0;\theta)
	\]
\end{definition}

The likelihood is not a probability nor probability density, but in general a non-negative function that indicates the "plausibility" (well, \emph{likelihood}) of extracting whichever $\theta$ given the measure $x_0$.
\todo{In the next example we are going to talk about distributions, such as the Bernoulli, we have not introduced yet. We should find a pleasing way to introduce them before.}
\begin{example}
	We want to determine the likelihood function for a binomial distribution. The binomial distribution gives the probability of observing $k$ successes in $n$ independent trials:
	\[
	\mathcal{B}=\binom{n}{k}p^k(1-p)^{n-k}.
	\]
	Applying the definition of Likelihood, we obtain:
	\[
	L_{x_0}(\theta)=p(x_0;\theta)\ \Rightarrow\ L_{n_0,k_0}(p)=\mathcal{B}(n_0,k_0;p).
	\]
	We consider three different cases:
	\begin{align*}
		L_{n=1,k=1}(p) &= p, \\
		L_{n=1,k=0}(p) &= 1-p, \\
		L_{n=50,k=25}(p) &= \binom{50}{25}p^{25}(1-p)^{25}.
	\end{align*}
	We plot these three likelihood functions in the graph below:
	\begin{figure}[H]
		\centering
		\input{images/binomial_likelihood.pgf}
		\caption{Likelihood functions for different values of $n$ and $k$.}
		\label{fig:binomial_likelihood}
	\end{figure}
	We can see that the likelihood suggests that the most likely values of $p$ are $1$, $0$, and $1/2$ for the three cases, respectively. This is consistent with intuition: if we observe one success in one trial, a naive estimate of the success probability is $p=1$.
\end{example}

\begin{example}\label{ex:uniform_likelihood}
	We want to determine the likelihood function for a uniform distribution. Consider a random variable distributed as:
	\[
	X \sim \mathcal{U}([0,m])
	\]
	whose probability density function is
	\[
	f(x;m)=
	\begin{cases}
		\frac{1}{m} & \text{if } x \in [0,m] \\
		0 & \text{otherwise}
	\end{cases}
	\]
	
	Given a single observation $x_0$, the likelihood function is obtained by viewing the density as a function of the parameter $m$:
	\[
	L_{x_0}(m) = f(x_0;m)
	\]
	
	From the definition of the density, we have:
	\[
	L_{x_0}(m)=
	\begin{cases}
		\frac{1}{m} & \text{if } m \ge x_0 \\
		0 & \text{if } m < x_0
	\end{cases}
	\]
	
	Equivalently, we can write it using the indicator function:
	\[
	L_{x_0}(m)=\frac{1}{m}\,\mathscr{F}(m>x_0)
	\]
	
	We observe that the likelihood is a decreasing function of $m$ on the interval $[x_0,+\infty)$, and is zero for $m < x_0$.	
	\begin{figure}[H]
		\centering
		\input{images/uniform_likelihood.pgf}
		\caption{Likelihood function for a single observation from a uniform distribution.}
		\label{fig:uniform_likelihood}
	\end{figure}
\end{example}\todo{I totally written this part from scratch, especially the conclusions. We should review this example in the future.}

\begin{example}
Likelihood Gauss
\end{example}

A few key properties of the Likelihood:
\begin{enumerate}
	\item $L$ only inherits the units from the observable, and can otherwise be discrete, continuous, multi-dimensional, independently from the observable's characteristics.
	\item The Likelihood does not have to be integrable nor normalized, and its integral has no interesting interpretation.
	\item When transforming variables, the Likelihood transforms like the pdf (so with $||J(x)||$), and therefore the ratio $\frac{L\left(\theta_1\right)}{L\left(\theta_2\right)}$ of two Likelihoods remains constant.
	\item Given a model for the observables, the functional form of the Likelihood is uniquely determined up to a multiplicative constant.
	\item If we apply an invertible transformation in the 	\emph{parameter} space, the Likelihood changes functional form but maintains its value ($L_x(f(\theta)) = L_x(\theta)$). 
	\begin{example}
		$L_{t}(\lambda)=\lambda e^{-\lambda t}=\lambda(\tau) e^{-\lambda(\tau) t}=\frac{e^{-t/\tau}}{\tau}\equiv L_{t}(\tau)$
	\end{example}
	In particular, the value of the parameter for which the Likelihood is maximum is invariant. 
	\item The Likelihood from the measure of many independent observables $x,y,\dots$ (or the same one measured many times independently), dependent on a certain set of parameters is the product of the individual Likelihoods. 
	\begin{example}
		Consider $x_i \sim \mathcal{U}([0,m])$ then:
		\[
		L_{\vec{x}}(m)=\frac{1}{m^N}\,\mathscr{F}(m>\max\{x_i\})
		\]
	\end{example}
	This means that more measurements can only better your Likelihood and that you can combine it with others to keep track of more observations. 
	\item $\ell =\ln{L}$ maintains all the same, nice properties as $L$, swapping the multiplicative constant for an additive one. For this reason, it's often used in calculations.
\end{enumerate}

\begin{example}
	Here he combines two Gaussian Likelihoods, see if you want.
\end{example}\todo{this}

It should be pretty clear now why the Likelihood is a natural tool if our goal is to draw conclusions about the parameters that govern the distributions from which we extract our observations. Though not interpretable as a probability, the likelihood still offers crucial information, and higher values can be interpreted as more plausible candidates for the parameter. Most importantly, where $L=0$, that value can be completely excluded as a possibility. Also, the multiplicative nature of repeated measures causes thinner and sharper peaks around the most likely values.

In general, we have a promising way to write a functional inference, where the statistic is the likelihood function itself. 

\begin{definition}[Likelihood principle]
	Central to the Likelihoodist perspective is the idea that the Likelihood function is the sole vehicle for evidence, leading to the following formal requirements for any rational inference: 
	\begin{enumerate}
		\item Any possible inference about $\theta$ derived from the observation $x_0$ can be based exclusively on $L(\theta)$.
		\item If two experiments yield the same $L(\theta)$ for a given $x$, the inferences that can be drawn about $\theta$ are identical.
		\item Any conclusion about $\theta$ that depends on anything other than $L(\theta)$ is inadmissible.
		\item Any conclusion about $\theta$ that depends on anything other than $L(\theta)$ must be rejected as mere irrationality or superstition.
	\end{enumerate}
\end{definition}

As we can tell, there are strong opinions on the matter and the argument was controversial, with the debate not yet resolved to this day. As far as we are concerned, we will take a much more neutral approach to the matter, integrating the concept of Likelihood with other techniques, adopting its helpful properties but violating the \emph{Likelihood principle}.

\subsection{Inferential use of the Bayes theorem}
\label{subsec:Inf_Bayes}

Let's assume, to start, that $\theta$ admits a probability distribution in $M$, namely $P(\theta)$, on which we hope to make some inferences. We write then, in this case,
	\[
	L_{x_0}(\theta)=p(x_0|\theta)
	\]
and using Bayes' theorem (\ref{th:Bayes}) we can write:
	\begin{equation}
	p(\theta|x_0) = \frac{\overbrace{p(x_0|\theta)}^{L_{x_0}(\theta)}p(\theta)}{\underbrace{p(x_0)}_{\int p(x_0|\theta)p(\theta)\dd\theta}}
	\label{def:posterior}
	\end{equation}
 
 In words: the \textbf{posterior} probability $p(\theta|x_0)$ is the product of the Likelihood times the \textbf{prior} probability $p(\theta)$ (times a multiplicative normalization factor $p(x_0$), which is the marginal probability of observing $x$ under our whole prior model). When it exists, it is naturally a better tool than the likelihood function for functional inferences as it builds upon it by adding the information tied to the distribution of $\theta$ itself. 
 
 While the Likelihood function is not necessarily integrable, the product $p(\theta)\cdot L_{x_0}(\theta)$ must be, as the posterior must be normalized. 
 
 
 If I can calculate this, i.e if $\theta$ admits a pdf and I am lucky enough to know it, then I know everything I will ever need to know and can calculate anything I may be interested in, like the \emph{probability}\sidenote{Not just "plausibility" or likelihood, the actual Probability} of any interval of $M$ or which is the $\theta$ that maximizes the posterior, for example. 
 
 \begin{example}
 	COVID
 \end{example}\todo{do}
 
 Often times, to avoid calculating the normalization, we talk about the ratio
 	\[
 	\frac{p(\theta_1|x_0)}{p(\theta_2|x_0)} = \overbrace{\frac{L_{x_0}(\theta_1)}{L_{x_0}(\theta_2)}}^{LR_{x_0}(\theta_1,\theta_2)} \cdot \frac{p(\theta_1)}{p(\theta_2)}
 	\]
 where $LR$ (\emph{Bayes' factor}) "updates" the relative probability between two hypotheses (different $\theta$) at each observation. Given the multiplicative nature of the Likelihood, this formula can be used to track the "evolution" of the relative probability. 
 
 \todo{Following example and also something about relative abundance?}
 \begin{example}
 	PID muoni e pioni etc...
 \end{example}

\subsubsection{Subjectivist Bayesian Inference}

Let's now talk about the more realistic scenario where $\theta$ does not admit a probability function (or it does, but we don't know it). The Bayesian (subjectivist) approach proceeds by \emph{assigning} a prior distribution to $\theta$, based on the \emph{degree of belief} $\pi(\theta)$. Then, the calculations proceed as before, and by treating the degree of belief as a probability we arrive to the posterior degree of belief:
\begin{equation}
	\pi_{x_0}(\theta) = \frac{L_{x_0}(\theta)\pi(\theta)}{\underbrace{\pi(x_0)}_{\sum_i p(x;\theta_i)\pi(\theta_i)}}
\label{eq:posterior_dob}
\end{equation}
In this expression, we now talk about the "credibility" of a certain $\theta$, remembering that the degree of belief is not a probability. Often times, it can be difficult to choose the right prior, and not always will it even be integrable (when it is not we call it an "improper prior"), though its product with the likelihood must by all means be integrable. 


Once again, we can talk about credibility ratios
 	\[
	\frac{\pi(\theta_1|x_0)}{\pi(\theta_1|x_0)} = \overbrace{\frac{L_{x_0}(\theta_1)}{L_{x_0}(\theta_2)}}^{LR_{x_0}(\theta_1,\theta_2)} \cdot \frac{\pi(\theta_1)}{\pi(\theta_2)}
	\]
 and the quantity $\frac{L_{x_0}(\theta_1)}{L_{x_0}(\theta_2)}$ is called \emph{belief-updating ratio} while $\frac{\pi(\theta_1)}{\pi(\theta_2)}$ hold the significance of \emph{betting odds} in the game theory world.
 
 This formulation is not the only way to handle statistical inferences when $\theta$ does not have a well known and well defined distribution, and we will tackle different and more complex methods in the following sections.\todo{Could be expanded on further, this whole subsection I mean not just this last sentence}

\section{Sufficient Statistics}

We have already talked about statistics, which are functions of random variables that do not depend on other parameters or factors that are not the observables themselves. We will now dive deeper into distinguishing the different types of statistics. 

\begin{definition}[Sufficient statistic]
	Let $p(x;\theta)$ be the pdf of $x$, depending on parameter $\theta$. Let $s(x)$ be a statistic (with dimensionality less than or equal to that of $x$) such that the probability of $x$ \emph{given} $s$ does not depend on $\theta$:
	\[
	p(x|s;\theta_1) = p(x|s;\theta_2) \forall\theta_1,\theta_2
	\]
	If this is the case, we say that the statistic $s$ is \emph{sufficient for} $\theta$.
\end{definition}
\begin{corollary}
	Every invertible function f(s) of a sufficient statistic is also sufficient, because $p(x|f) = p(x | f(s)) = p(x|s)$\sidenotemark
\end{corollary}
	\sidenotetext{Since $f(s)$ is an invertible function, it uniquely determines $s$, ergo: conditioning on $f(s)$ is the same as conditioning on $s$}
\begin{definition}[Ancillary statistic]
	The less used complementary cousin of the sufficient statistic: if $s(x)$ is a statistic such that its pdf is not dependent on $\theta$:
	\[
	p(s;\theta_1) = p(s;\theta_2) \, \forall \theta_1,\theta_2
	\]
	Then we say that the statistic is \emph{ancillary} for $\theta$
\end{definition}
\begin{theorem}[Factorization theorem]
	The statistic $s(x)$ is sufficient for $\theta$ if and only if $p(x;\theta)$ can be factorized in a part that does not depend on $\theta$ and a part that depends on the observable $x$ only through $s(x)$ and $\theta$. Synthesized:
	\[
	s \text{ sufficient}\iff p(x;\theta) = f(x)\cdot g(s(x),\theta) \quad \text{with } f,g\geq0
	\]
	Important note: the domain of $x$ must be \emph{independent of} $\theta$ (see \ref{ex:uniform_sufficient_statistics})
\label{th:factorization}
\end{theorem}

\begin{example}[Proof that factorization implies sufficiency, factorization theorem $\Leftarrow$]
	\begin{align*}
	p(s;\theta) &= \int_X p(s,x;\theta)\dd x = \int_X p(s|x;\theta)p(x;\theta)\dd x \\
	&= \sidenotemark \int_X \delta(s-s(x))p(x;\theta) \dd x = \int_X \delta(s-s(x))f(x)g(s(x),\theta) \dd x \\
	&= g(s,\theta)\int_X \delta(s-s(x))f(x)\dd x
	\end{align*}
	 
	\begin{align*}
	p(x|s;\theta) &= \frac{p(s|x;\theta)\cdot p(x;\theta)}{p(s;\theta)} = \frac{\delta(s-s(x))\cdot p(x;\theta)}{p(s;\theta)} =\\
	&= \frac{\delta(s-s(x))\cdot f(x)\cdot g(s,\theta)}{g(s,\theta)\int_X \delta(s-s(x))f(x)\dd x} = \frac{\delta(s-s(x))\cdot f(x)}{\int_X\delta(s-s(x))f(x)\dd x}	
	\end{align*}
	which is independent of $\theta$. \qed	
\end{example}
\sidenotetext{Since $s(x)$ is deterministic, the conditional distribution of $s$ given $x$ is a Dirac delta.}

\begin{example}[Proof that sufficiency implies factorization, factorization theorem $\Rightarrow$]
	If
	\[
	p(x | s;\theta_1) = p(x | s;\theta_2) \, \forall \theta_1,\theta_2
	\]
	then 
	
	\[
	p(x;\theta)=\frac{p(x|s;\theta)\cdot p(s;\theta)}{p(s|x;\theta)} =\sidenotemark \frac{p(x|s)}{p(s|x)}\cdot p(s;\theta)
	\]
	\qed
\end{example}
\sidenotetext{$p(s|x; \theta)=\delta(s-s(x))$ as said before and $p(x|s;\theta)$ is independent of $\theta$ by hypothesis}

In summary, sufficient statistics "compress" the data without losing any information on the parameter $\theta$. Once I have a sufficient statistic of the data I can make equally valid inferences on $\theta$ using only the statistic as I would with the whole dataset, and therefore have no reason to keep the whole set, if my goal is inferring $\theta$. 

\begin{example}\label{ex:uniform_sufficient_statistics}
	Consider $x$ following a uniform distribution $\mathcal{U}([0,m])$.
	
	\textbf{A misleading approach.} One might naively write the pdf as
	\[
	p(x;m)=\frac{1}{m}, 
	\quad
	p(\vec{x};m)=\frac{1}{m^N},
	\]
	and conclude that the likelihood is constant in $\vec{x}$. This could suggest that any statistic (for instance $s=x_1$) is sufficient.
	
	However, this reasoning is incorrect because it neglects the support of the distribution, which depends on $m$.

	\textbf{Correct formulation.} The pdf must include indicator functions:
	\[
	p(x;m)=\frac{1}{m}\,\mathbb{I}_{[0,m]}(x),
	\quad
	p(\vec{x};m)=\frac{1}{m^N}\prod_{i=1}^N \mathbb{I}_{[0,m]}(x_i).
	\]
	
	We can rewrite the joint indicator as
	\[
	\prod_{i=1}^N \mathbb{I}_{[0,m]}(x_i)
	=
	\mathbb{I}_{\{0 \le \min_i x_i\}}
	\mathbb{I}_{\{\max_i x_i \le m\}}.
	\]
	
	Thus,
	\[
	p(\vec{x};m)
	=
	\frac{1}{m^N}
	\mathbb{I}_{\{0 \le \min_i x_i\}}
	\mathbb{I}_{\{\max_i x_i \le m\}}.
	\]

	We can now write
	\[
	p(\vec{x};m)
	=
	\underbrace{\mathbb{I}_{\{0 \le \min_i x_i\}}}_{f(\vec{x})}
	\cdot
	\underbrace{\frac{1}{m^N}\mathbb{I}_{\{\max_i x_i \le m\}}}_{g(\max\{\vec{x}\},m)}.
	\]
	
	Using the factorization theorem we can see that $\max\{\vec{x}\}$ is a sufficient statistic for $m$. We also note that clearly $g(\max\{\vec{x}\},m)$ is not the pdf of $\max\{\vec{x}\}$ given $m$.
\end{example}
\todo{I'm still a bit confused by the requirement of a domain of $x$ independent of $\theta$. When presenting the theorem that is said to be a necessary condition, but in the example above we apply the theorem anyway and it is said to work. In the lecture notes the focus of this requirement shifts from a condition necessary for the theorem to work to a condition for the $g$ to be unique and a proper pdf. This is a bit FRANK to me.}

The following principle sums up the concept quite nicely:

\begin{theorem}[Sufficiency principle]
	Let $s(x)$ be sufficient for $\theta$. Every possible inference on $\theta$ deriving from $x_0$ must be based exclusively on $s(x_0)$
\end{theorem}

Interesting to note that the Likelihood principle implies the sufficiency principle, if we let $L_x(\theta)=f(x)\cdot g(s(x),\theta)$. If we know $s(x)$ for all values of $x$, then we know $L_x(\theta)$ (up to an irrelevant multiplicative constant) for all values of $x$, which is nice in general but weaker than the Likelihood principle where we only need to know $p(x;s)$ for one observed value $x_0$, not all $x$. 

 It's hard to overstate the simplicity yet power of this statement, and in the course of this text we will find ourselves dealing with sufficient statistics more often than entire lengths of data. 

\begin{definition}[Minimal sufficient statistic]
	A sufficient statistic is said to be \emph{minimal} if:
	\[
	\forall s_i:s_i\text{ sufficient}\Rightarrow s=s(s_i)
	\]
	This is equivalent to asking that the information can't be further reduced or compressed. 
\end{definition}
\todo{There will be exercises that clear these up quite nicely, add them here or ref them}

\todo{Conditionality principle, Birnbaum argumentation, both things I'd happily skip... Shall we?}

So now we have talked at lengths about this beautiful concept called the sufficient statistic and how it is going to be the resolution to all of our problems, and provided some examples of well-known sufficient statistics for important distributions.\todo{Have we? Either add them or change this phrase to reference the appendix where all the important distributions and their properties will be.} All we need now is a way to more generally find the sufficient statistic for a given distribution, and luckily there is a way through the following theorem:

\begin{theorem}[Darmois theorem]
	Given a real observable $x$ distributed according to $p(x;\theta)$ with domain independent from $\theta$, we suppose that, given $N$ observations ${x_k}_N$, there exists a sufficient statistic for $\theta$ $s_i\left({x_k}_N\right)$ with values in $\R^n$ (with $n$ independent from $N$).
	
	In this case, the pdf can be written as:
	\[
	p(x;\theta) = \exp{\left[ \sum_{i=1}^n \alpha_i(x)a_i(\theta) + \beta(x) +\gamma(\theta)\right]}
	\]
	
	And the sufficient statistic $s$ can be expressed as
	\[
	s_i(x)=\sum_{k=1}^N \alpha_i(x_k)
	\]
\label{th:darmois}
\end{theorem}

Basically, what this theorem is saying is that if I can write my pdf in exponential form, and seperate it into a sum terms as shown in \ref{th:darmois} then I can identify the sufficient statistic for $p$. This is great and will be used when possible, however it's easy to imagine that the class of functions that can be written as exponentials is not as large as we'd like.

\section{Fisher Information}

Let's take a step back and remember what we are here for. In this chapter, we want to explore different functions that apply a "plausibility score" to different states of a certain parameter. We saw the Likelihood and digressed into talking about sufficient statistics. Let's put these together and explore a new concept in the field of functional inferences:

\begin{definition}[Fisher score function]
	\[
	S_x(\theta) = \frac{\partial \ln L_x(\theta)}{\partial \theta} = \frac{\partial \ln p(x;\theta)}{\partial \theta}
	\]
\label{def:Fisher_S}
\end{definition}

This expression is a function of $\theta$, though for fixed $\theta$ it can be by all means considered a statistic, and for example, at a fixed $\theta$:
	\[
	\E\left[ S_x(\theta) \right] = \int \frac{\partial \ln p(x;\theta)}{\partial \theta}p(x;\theta)\dd x
	\]

We note that, if the domain of $x$ is independent of $\theta$ and we have sufficient regularity in $p$\sidenote{To be precise, the regularity conditions are as follows:\begin{itemize}
		\item The partial derivative of $p(x; \theta)$ with respect to $\theta$ exists almost everywhere (It can fail to exist on a null set, as long as this set does not depend on $\theta$)
		\item The integral of $p(x; \theta)$ can be differentiated under the integral sign with respect to $\theta$
		\item \textbf{The support of $p(x; \theta)$ does not depend on $\theta$}
\end{itemize}} then we can commute the derivative with the integral and we get

	\begin{align*}
	\E\left[ S_x(\theta) \right] &= \int\frac{1}{p(x;\theta)}\frac{\partial p(x;\theta)}{\partial \theta} p(x;\theta) \dd x \\
	&= \int \frac{\partial p(x;\theta)}{\partial\theta} \dd x = \frac{\partial}{\partial \theta}\underbrace{\int p(x;\theta)\dd x}_1 \equiv 0\quad \forall \theta
	\end{align*}

This means that under sufficient regularity (and if the domain of the observable is independent of the parameter), then the expected value of the Fisher score is identically null.

\begin{definition}[Fisher Information function]\label{def:fisher_information}
	\[
	I(\theta) = \E\left[ S^2_x(\theta) \right] = \E\left[\left( \frac{\partial \ln L_x(\theta)}{\partial\theta}\right)^2\right]
	\]
\end{definition}

Under the same principles of sufficient regularity, given that the expected value of the Fisher score is zero, then \textbf{the Fisher Information will be equal to the variance of the score}. 

	\[
	\Var\left[S_x(\theta)\right] = I(\theta)
	\]

Generalizing to more parameters, we have the \emph{Fisher Information matrix} defined as:
	\[
	I_{ij}(\theta) = \E\left[ \frac{\partial \ln L_x(\theta)}{\partial \theta_i}\cdot \frac{\partial \ln L_x(\theta)}{\partial \theta_j}\right]
	\]

This matrix, under the usual regularity requisites, can be expanded in the following way:

\begin{align*}
	I_{ij}(\theta) &= \int \frac{\partial \ln p(x;\theta)}{\partial \theta_i} \frac{\partial \ln p(x;\theta)}{\partial \theta_j} p(x;\theta)\dd x = \int \frac{\partial \ln p(x;\theta)}{\partial \theta_i} \frac{\partial p(x;\theta)}{\partial \theta_j} \dd x \\
	&= \int\left[ \frac{\partial}{\partial\theta_i}\left(p(x;\theta)\frac{\partial \ln p(x;\theta)}{\partial \theta_i}\right) - p(x;\theta)\frac{\partial^2\ln p(x;\theta)}{\partial \theta_i \partial\theta_j}\right]\dd x \\
	&= \frac{\partial}{\partial \theta_i}\E\left[S_x(\theta)\right] - \E\left[\frac{\partial^2 \ln p(x;\theta)}{\partial \theta_i\partial\theta_j}\right] \\ 
	&= \E\left[-\frac{\partial^2\ln L(\theta)}{\partial\theta_i\partial\theta_j}\right]
\end{align*}

That, in the usual case of a single parameter, is:
\[
	I(\theta)=\E\left[-\frac{\partial^2\ln L(\theta)}{\partial\theta^2}\right]
\]

This expression makes it clear that the Fisher Information is linear in $\ln L$, and since the Likelihood is multiplicative in the number of observations, the Fisher Information is additive (\textbf{under the usual assumptions}). Since we want the Information to represent how much the observations tell us about the parameter $\theta$, this is a property that we happily accept.\todo{We should review this section. I think there's margin of improvement and we could make the "Fisher Information is additive" part more clear and rigorous.}

To put everything together, let's extend the concept of Fisher Information to \emph{statistics}, given a statistic $s(x)$ let's call the \emph{Information} of $s$ on $\theta$ the quantity:
	\[
	I_s(\theta) = \E \left[\left( \frac{\partial\ln p(s(x);\theta)}{\partial \theta}\right)^2 \right] 
	\]
	
If we consider a sufficient statistic for $\theta$, then we have:

\begin{align*}
	p(x;\theta)&=p(x|s;\theta)p_s(s;\theta)=p(x|s)p_s(s;\theta) \\ &\Rightarrow \ln p(x;\theta) = \ln p_s (s;\theta) + \ln p(x|s) \\
	&\Rightarrow \frac{\partial}{\partial \theta}\ln p(x;\theta) = \frac{\partial}{\partial\theta}\ln p_s(s,\theta) + \cancelto{0}{\frac{\partial}{\partial \theta}\ln p(x|s)} = \frac{\partial}{\partial\theta}\ln p_s(s;\theta)\\
	&\Rightarrow \E\left[\left(\frac{\partial\ln p(x;\theta)}{\partial\theta}\right)^2\right] = \E\left[\left(\frac{\partial\ln p_s(s;\theta)}{\partial\theta}\right)^2\right]\\
	&\Rightarrow I_s(\theta) = I(\theta)
\end{align*}

Therefore, the Information of a sufficient statistic is the same as the Information of the original data itself. Given a generic statistic (not necessarily sufficient), it can be shown that
	\[
	I_{s(x)}(\theta)\leq I(\theta)
	\]
and if the statistic $s(x)$ is an invertible function of $x$, we have:
	\[
	I(\theta)\equiv I_{x(s)}\leq I_s(\theta)
	\]
which tells us that:
	\[
	I(\theta) = I_s(\theta)
	\]
Thus implying that an invertible transformation does not change the Information.

All of these are properties that we would expect of a function that symbolizes the information of an observable/statistic, and if I cannot find a sufficient statistic for my observable a good idea might be to find a statistic that minimizes the Fisher Information loss.

\subsubsection{Asymptotic distribution}

While we have established that the Fisher score is a statistic with zero mean and variance equal to the Fisher Information, its true power in inference is revealed as the sample size $n$ becomes large. Since the log-likelihood of a sample of $n$ independent and identically distributed (i.i.d.) observations is the sum of the individual log-likelihoods, the Fisher score $S_n(\theta)$ is likewise a sum of i.i.d. random variables:

	\[
	S_n(\theta) = \sum_{i=1}^n \frac{\partial \ln p(x_i; \theta)}{\partial \theta}
	\]

By invoking the Central Limit Theorem, we can characterize the limiting behavior of this sum. For a fixed $\theta$, as $n \to \infty$, the distribution of the Score function scaled by $1/\sqrt{n}$ converges to a normal distribution with variance equal to the unit Fisher Information. In practical terms, for a large sample, we say that the total score is asymptotically normal:

\begin{equation}
	S_n(\theta) \sim \N\left(0, I_n(\theta)\right)
	\label{eq:Score_asymp}
\end{equation}

 This will be an important result that we use in the following chapters, and once again stabilizes why the 
 
 \vspace*{1cm}
 \begin{center}
 	\textbf{Exercises at: \ref{exer:functional_inf}}
 \end{center}